\subsection{原始射に支持を持つ変形による全次数の対称性保持}

四対象 $(0,1,3,4)$ で、原始射 $w_0:0\to3$ と次数一射の積だけを動かす
部分空間 $W$ は五十四次元、その出力を純粋な四次多項式に制限した $W_0$ は
四十五次元である。
三重合成 $0\to1\to3\to4$ の最初の積には $w_0$ が現れないため、
全 $F,G\in W$ について $d^2F=0$、$F\circ G=0$ である。
この支持の判定は、有限次数の検証を全次数の証明に変えるために再利用できる。

不変cocycle $U_i$ と導分補正 $E_i$ について $d^1E_i=[\delta,U_i]$ とする。
$E_i$ が $w$ を含む射を純粋多項式へ送るだけなら、
全 $V\in W_0$ に対して $[E_i,V]=0$ である。
$B_{0,3}$ の四本の不変方向は全て $W_0$ に入り、
$\mu+\sum u_iU_i$ と $\delta+\sum u_iE_i$ がそのまま全次数の族となる。
接空間次元と一致するため、対称性保持部分全体は滑らかな四変数形式空間である。

$B_{1,1}$ の $\ker L\cap\operatorname{im}L$ の五本では、
二次の導分誤差は $W_0$ に入り、階数三十八の
$L_\delta:W_0\to W_0$ の像に入る。
有理数上で $L_\delta V_{ij}=-([E_i,U_j]+[E_j,U_i])$ を解く
（対角項は一項だけ）と、積の二次補正だけで全次数の族を得る。
三次以上の導分条件と全結合則は支持から零である。
この族は滑らかな五次元閉形式部分を与え、対称性保持部分の次元が正確に五と決まる。
両族の導分は $w$ フィルトレーションを保ち、その補正は次数を下げるため、
一様に十二乗が零であり、実際の $\mathbf G_a$ 作用を持つ。

五次元の滑らかな部分を $t=0$ と直線化すれば、全跳躍環は
\[
 A[[t]]/(t(t-b),tJ)
   =A\oplus(A/J)t,\qquad
 A=\mathbf C[[u_1,\ldots,u_5]],\quad
 b\in\mathfrak m_A^2,\quad J\subset\mathfrak m_A^2
\]
の形になる。証明は、定義イデアルが $(t)$ に含まれることと、
二次項 $t^2$ を持つ関係を $t$ で割った元への形式的陰関数定理である。
この有限表示は、横断族を構成する前の候補を限定するために有用である。
v39では横断方向を含む全次数の有理式族により $b=J=0$ と決まった。
従って実際の環は $A[\varepsilon]/(\varepsilon^2)$ であり、$A$ 上自由階数二である。
